\section{把这一切统一起来}

\subsection{一个稳定的工程模式}

KenLM、fastText、DSIR、Bloom filter、MinHash、LSH，看起来像不同主题，实际上都在复用同一类工程模板：先用便宜的统计量压缩问题，再把大规模搜索变成分桶、排序、阈值判断或者重采样。

\begin{figure}[H]
\centering
\begin{tikzpicture}[node distance=1.2cm, every node/.style={font=\small, align=center}]
\node[draw, rounded corners, fill=green!10, minimum width=3cm, minimum height=0.85cm] (fit) {fit score model};
\node[draw, rounded corners, fill=yellow!16, below=of fit, minimum width=3cm, minimum height=0.85cm] (score) {score raw data};
\node[draw, rounded corners, fill=blue!14, below=of score, minimum width=3cm, minimum height=0.85cm] (act) {keep / resample / bucket};
\node[draw, rounded corners, fill=red!10, below=of act, minimum width=3cm, minimum height=0.85cm] (done) {better corpus};
\draw[-{Latex[length=3mm]}, thick] (fit) -- (score);
\draw[-{Latex[length=3mm]}, thick] (score) -- (act);
\draw[-{Latex[length=3mm]}, thick] (act) -- (done);
\end{tikzpicture}
\caption{统一框架：先训练一个粗糙但快的打分器，再把它应用到大规模原始数据上}
\end{figure}

\begin{table}[H]
\centering
\small
\renewcommand{\arraystretch}{1.15}
\begin{tabular}{p{2.2cm}p{4.1cm}p{4.0cm}p{3.3cm}}
\toprule
\textbf{方法} & \textbf{score} & \textbf{action} & \textbf{典型用途} \\
\midrule
KenLM & $p_T(x)$ / perplexity & threshold / sort & language / quality filtering \\
fastText & $p(T\mid x)$ & threshold / classify & language / toxicity / quality \\
DSIR & $p_T(x)/p_R(x)$ & weighted resampling & more principled data selection \\
Bloom filter & membership bits & approximate reject & exact dedup \\
MinHash + LSH & collision probability & candidate recall & near dedup \\
\bottomrule
\end{tabular}
\caption{同一套机制在不同任务里的落点}
\end{table}

\begin{warningbox}{最重要的现实约束}
这些方法不是在追求“数学上最优”，而是在追求“算得起、筛得动、错得起”。大规模数据处理的核心，就是在 compute、memory、误杀率之间做足够清楚的取舍。
\end{warningbox}

